%MIT OpenCourseWare: https://ocw.mit.edu
%18.100A / 18.1001 Real Analysis, Fall 2020
%License: Creative Commons BY-NC-SA 
%For information about citing these materials or our Terms of Use, visit: https://ocw.mit.edu/terms.

\subsection*{Sequences of Function}

\noindent\underline{\textbf{Power Series}}

\begin{remark}
Power series motivate the general discussion of sequences of functions.
\end{remark}

\begin{definition}[Power series]
A power series about $x_0$ is a series of the form 
\[
\sum_{m=0}^\infty a_m (x-x_0)^m.
\]
\end{definition}

\begin{theorem} 
Suppose
\[
R = \lim_{m\to \infty} |a_{m}|^{1/n}
\]
exists, and let 
\[
p = \begin{cases}
\frac{1}{R} & R>0 \\
\infty & R = 0.
\end{cases}
\]
Then, $\sum a_m(x-x_0)^m$ converges absolutely if $|x-x_0|<p$ and diverges if $|x-x_0|>p$.
\end{theorem}

\begin{definition}[Radius of Convergence]
In the above theorem, we define $p$ to be the \vocab{radius of convergence}.
\end{definition}
\textbf{Proof}: We have 
\[
\lim_{n\to \infty}|a_m(x-x_0)^m|^{1/m} = R|x-x_0|,
\]
and the theorem follows by the Root test. \qed 

Suppose $\sum a_m(x-x_0)^m$ is a power series with radius of convergence $p$. Furthermore, define \\$f:(x_0-p,x_0+p) \to \RR$ such that 
\[
f(x) := \sum_{m=0}^\infty a_m (x-x_0)^m.
\]
Then, $f$ is a limit of a sequence of functions 
\begin{align*}
    f(x) &= \lim_{n\to \infty} f_n(x),
\intertext{for $x\in (x_0-p,x_0+p)$ and where}
    f_n(x) &= \sum_{m=0}^n a_m(x-x_0)^m.
\end{align*}

\begin{example}
For example, we have 
\[
f(x) = \frac{1}{1-x} = \sum_{m=0}^\infty x^m.
\]
\end{example}

\begin{question}
This concept begs a number of questions:
\begin{enumerate}
    \item Is $f$ continuous? 
    \item Is $f$ differentiable, and does $f' = \lim_{n\to \infty} f_n'$? 
    \item If 1. is true, does 
    \[
    \int_a^b f = \lim_{n\to \infty}\int_a^b f_n?
    \]
\end{enumerate}
\end{question}

\noindent These questions will be the key motivator for the last section of this course.

\noindent\underline{\textbf{Pointwise and Uniform Convergence}}

We now consider a setting far more general than power series.

\begin{definition}[Pointwise Convergence]
For $n\in \NN$, let $f_n: S\to \RR$. Let $f:S\to \RR$. We say that $\{f_n\}$ converges pointwise to $f$ if for all $x\in S$,
\[
\lim_{n\to \infty} f_n(x) = f(x).
\]
\end{definition}

\noindent Let's consider some examples.
\begin{enumerate}
    \item Let $f_n(x) = x^n$ on [0,1]. Then, 
    \[
    \lim_{n\to \infty} f_n(x) = \begin{cases}
    0 & x\in [0,1) \\
    1 & x=1
    \end{cases}.
    \]
    Thus, $\{f_n\}$ converges to the above pointwise function. Hence, notice that a sequence of continuous functions may not converge pointwise to a continuous function!
    \item Let $f_n(x) = \sum_{m=0}^n x^m$ for $x\in(-1,1)$.
    Then, 
    \[
    \lim_{n\to \infty} f_n(x) = \lim_{n\to \infty} \sum_{m=0}^n x^m = \frac{1}{1-x}.
    \]
    Hence, pointwise, this sequence converges to its power series (see the above example).
    \item Let $f_n(x):[0,1] \to \RR$ be defined by 
    \[
    f_n(x) =\begin{cases}
    4n^2x & x\in \left[0,\frac{1}{2n}\right] \\
    4n-4n^2x & x\in \left[\frac{1}{2n}, \frac{1}{n}\right] \\
    0 & x\in \left[\frac{1}{n}, 1\right]
    \end{cases}.
    \]
    We can picture this sequence (on the next page)
    \begin{figure}[h]
    \centering
    \includegraphics{pics/fig12.PNG}
\end{figure}

Then, $\lim_{n\to \infty} f_n(0) = \lim_{n\to \infty} 0 = 0$. Let $x\in (0,1]$. Let $N\in \NN$ such that $\frac{1}{N}< x$. Then, for all $n\geq N$, 
\[
f_n(x) = 0.
\]
Therefore, 
\[
\{f_n(x)\} = f_1(x), \dots, f_{N-1}(x), 0, 0, 0, \dots.
\]
Hence, $\lim_{n\to \infty}f_n(x) = 0$ for all $x\in [0,1]$. Thus, $\{f_n\}$ converges pointwise to $f(x) = 0$ on $[0,1].$
\end{enumerate}

\begin{definition}[Uniform Convergence]
For $n\in \NN$, let $f_n:S\to \RR$, and let $f:S\to \RR$. Then, we say $f_n$ \vocab{converges to $f$ uniformly} or \textbf{converges uniformly to $f$} if $\forall \epsilon>0$ $\exists M\in \NN$ such that for all $n\geq M$ $\forall x\in S$,
\[
|f_n(x) - f(x)|<\epsilon
\]
\end{definition}

\begin{theorem}
If $f_n:S\to \RR$, $f:S\to \RR$, and $f_n\to f$ uniformly, then $f_n\to f$ pointwise.
\end{theorem}
\textbf{Proof}: Let $c\in S$ and let $\epsilon>0$. Then, $f_n\to f$ uniformly implies that there exists $M_0\in \NN$ such that for all $n\geq M, \forall x\in S$, $|f_n(x)-f(x)|<\epsilon$. Choose $M=M_0$. Then, $\forall n\geq M$, 
\[
|f_n(c) - f(c)|< \epsilon.
\]
Thus, $\lim_{n\to \infty}f_n(c) = f(c)$ for all $c\in S$, and therefore $f_n\to f$ pointwise. \qed